\chapter{Introduction}

Text-to-image diffusion models, Stable Diffusion v1.5 (SDv1.5)~\cite{Rombach_2022_CVPR}
in particular, have achieved widespread deployment while relying on massive web-scraped
corpora such as Re-LAION-5B~\cite{Schuhmann_2022_Laion5b}, whose automated construction
introduces systematic biases that reflect societal
inequities~\cite{Birhane_2021_MultiModalDatasetsMisogynyPornography}. Existing research
documents demographic bias in generative
outputs~\cite{Luccioni_2023_AnalyzingSocietalRepresentationsDiffusionModels,
    Bianchi_2023_TTIGenAmplifiesDemographicStereotypes} but treats models as black boxes,
ignoring the \emph{spurious features} — low-level non-semantic characteristics such as
colour statistics, texture patterns, edge distributions, frequency domain properties, and
compression artefacts — that correlate with demographic attributes and enable biased
shortcuts~\cite{Meister_2023_GenderArtifactsinVisualDatasets,
    Zeng_2024_UnderstandingBiasinLargeScaleVisualDatasets}. This thesis addresses these gaps
through a mechanistic investigation of spurious features in Re-LAION-5B and SDv1.5,
applying established methodologies to expose low-level features that distinguish dataset
sources and correlate with demographic attributes, then evaluating targeted
inference-time interventions without model retraining. The analysis reveals that
Re-LAION-5B and SDv1.5 encode demographic information through qualitatively different
mechanisms, with direct consequences for bias attribution and mitigation. Five objectives
structure the work: (\textit{O1})~dataset acquisition and annotation across 200
ISCO-08-aligned occupations; (\textit{O2})~spurious feature detection via classifier-based
transformation analysis; (\textit{O3})~demographic correlation assessment and
occupational representation comparison against real-world statistics;
(\textit{O4})~inference-time debiasing of SDv1.5; and (\textit{O5})~synthesis of how
spurious features and demographic biases are transformed by the SDv1.5 generative
process.


% OLD INTRO
% \chapter{Introduction}

% Text-to-image diffusion models, Stable Diffusion~\cite{Rombach_2022_CVPR} in
% particular, have achieved widespread deployment while relying on massive web-scraped
% corpora such as Re-LAION-5B~\cite{Schuhmann_2022_Laion5b}, whose automated construction
% introduces systematic biases that reflect societal
% inequities~\cite{Birhane_2021_MultiModalDatasetsMisogynyPornography}. Existing research
% documents demographic bias in generative
% outputs~\cite{Luccioni_2023_AnalyzingSocietalRepresentationsDiffusionModels,
% Bianchi_2023_TTIGenAmplifiesDemographicStereotypes} but treats models as black boxes,
% ignoring the \emph{spurious features} — low-level non-semantic characteristics such as
% colour statistics, texture, and edge distributions — that correlate with demographic
% attributes and enable biased
% shortcuts~\cite{Meister_2023_GenderArtifactsinVisualDatasets,
% Zeng_2024_UnderstandingBiasinLargeScaleVisualDatasets}. This thesis addresses this gaps
% through a mechanistic investigation of spurious features in Re-LAION-5B and SDv1.5,
% applying established methodologies to expose low-level features that distinguish dataset
% sources and correlate with demographic attributes, then evaluating targeted
% inference-time interventions without model retraining. Five objectives structure the
% work: (\textit{O1})~dataset acquisition and annotation across 200 ISCO-08-aligned
% occupations; (\textit{O2})~spurious feature detection via classifier-based
% transformation analysis; (\textit{O3})~demographic correlation assessment against
% real-world statistics; (\textit{O4})~inference-time debiasing of SDv1.5; and
% (\textit{O5})~synthesis of how training data biases propagate through the generative
% process.